\documentclass[11pt]{article}
\usepackage[T1]{fontenc}
\usepackage{lmodern,microtype,amsmath,amssymb,amsthm,mathtools,booktabs,array}
\usepackage[letterpaper,margin=1in]{geometry}
\usepackage{xurl}
\usepackage[colorlinks=true,allcolors=blue]{hyperref}
\newtheorem{theorem}{Theorem}[section]
\newtheorem{lemma}[theorem]{Lemma}
\newtheorem{proposition}[theorem]{Proposition}
\newcommand{\ceil}[1]{\left\lceil #1\right\rceil}
\newcommand{\ind}{\mathbf 1}
\DeclareMathOperator{\wt}{wt}
\title{Three finite lower bounds for superpermutations}
\author{Jay Pantone}
\date{October 2026}
\hypersetup{pdftitle={Three finite lower bounds for superpermutations},pdfauthor={Jay Pantone}}
\begin{document}
\maketitle
\begin{abstract}
We prove lower bounds of $4{,}033{,}379$, $6{,}746{,}627{,}228$, and
$93{,}891{,}141{,}150$ for superpermutations on ten, thirteen, and fourteen
symbols. The first two follow from finite exclusions in Halfar's model
of chain pieces, with a direct passage to unrestricted words. The third
uses the small amount of slack in our general lower bound to force all
runs into one insertion-loop partition, where their sizes are incompatible
with complete coverage. The ten- and fourteen-symbol proofs use Lean's
kernel throughout. The thirteen-symbol proof uses \nolinkurl{native_decide}
for 51 finite computations; the completeness argument and the reduction
to words are checked by the kernel.
\end{abstract}

\noindent\textbf{AI disclosure.} The research presented here, the Lean
formalization, and this document were generated entirely by AI systems
under Jay Pantone's direction. The inherited mathematical work and
software retain their original credits. Independent AI review is not
human peer review.

Only the ten- and fourteen-symbol results are submitted to Palomar. The
thirteen-symbol result is included in this paper for context; its native
computation assumptions are outside the submitted proof library.

\section{The three bounds}
A \emph{superpermutation} on $n$ symbols is a word that contains every
permutation of those symbols as a contiguous factor. Let $S(n)$ denote
the minimum length of such a word. Our general lower bound
$B(n)$ in~\cite{returns} comes from restrictions on the local structure
of a minimum overlap path. There are two ways to improve its value at
a fixed $n$. We can exclude more local configurations by a finite
calculation, or we can examine the configurations that would be necessary
for equality. We use both methods here.

\begin{theorem}\label{thm-main}
We have
\[
 S(10)\ge4{,}033{,}379,\qquad
 S(13)\ge6{,}746{,}627{,}228,\qquad
 S(14)\ge93{,}891{,}141{,}150.
\]
\end{theorem}

\begin{table}[ht]
\centering
\begin{tabular}{rrrr}
\toprule
$n$ & General bound $B(n)$ & Bound here & Increase\\
\midrule
10 & $4{,}033{,}374$ & $4{,}033{,}379$ & 5\\
13 & $6{,}746{,}626{,}531$ & $6{,}746{,}627{,}228$ & 697\\
14 & $93{,}891{,}141{,}149$ & $93{,}891{,}141{,}150$ & 1\\
\bottomrule
\end{tabular}
\caption{Comparison with the general theorem of~\cite{returns}.}
\end{table}

\enlargethispage{\baselineskip}
The finite search framework is due to Halfar~\cite{halfar}. His
all-length completeness theorem is essential: searching a collection of
short paths alone would not give a bound for arbitrary words. Our
contribution for $n=10$ and $n=13$ is a stronger parameter certificate
and its connection to the hereditary segment and forest argument of
our lower-bound framework. The fourteen-symbol argument is different.
Its last step uses the actual occupied rotation classes, information
that a list of numerical deficits does not retain.

Halfar's announcement of 6 October 2026~\cite{halfar-announcement}
gave $S(10)\ge4{,}033{,}378$ and $S(14)\ge93{,}891{,}141{,}008$.
The two bounds submitted here exceed those announced values by $1$
and $142$, respectively. This comparison is with that dated
announcement.

\section{From words to runs}
For permutations $u$ and $v$ of an $n$-symbol alphabet, the \emph{overlap
weight} $w(u,v)$ is $n$ minus the largest length of a suffix of $u$ that
is a prefix of $v$. An ordering of all permutations therefore gives a
superpermutation of length $n$ plus its total overlap weight. Conversely,
listing the first occurrences of the permutations in any superpermutation
gives an overlap path of no greater length. Thus
\[
 S(n)=n+\min_P\wt(P),
\]
where the minimum is over Hamiltonian paths in the overlap graph.

Weight-one steps rotate a permutation to the left. The $n$ rotations of
a permutation form a \emph{rotation class}. A traversal of all its members
is a \emph{block}. If the block is entered at
$(x_1,\ldots,x_{n-1},s)$, it is exited at
$(s,x_1,\ldots,x_{n-1})$. A weight-two step to a new rotation class must
then enter at
\[
 (x_2,\ldots,x_{n-1},x_1,s).
\]
The other possible completion of the overlap would revisit the old
rotation class. In particular, the final symbol $s$ is unchanged.

A maximal succession of blocks joined by weight-two steps is a
\emph{run}. It contains at most $n-1$ blocks. Its \emph{deficit} is
$n-1$ minus the number of its blocks; a run of deficit zero is
\emph{full}. A \emph{light segment} is a succession of runs with
weight-three joins. If it has $r$ runs and total deficit $D$, then it
contains $(n-1)r-D$ rotation classes.

An arbitrary overlap path need not consist of intact blocks. We use the
reduced minimum-path decomposition of~\cite{framework,returns}, which
extracts block-aligned components and accounts for the remaining
vertices. The part of that result needed below is as follows. Suppose
that every light segment, and every consecutive factor of such a
segment, satisfies
\begin{equation}\label{eq-affine}
 D\ge\delta r-\beta,
 \qquad 0<\delta<1,\quad 1\le\beta\le n-1-\delta.
\end{equation}
Put $\alpha=n-1-\delta$. There is an integer budget $b\ge0$ for each
reduced minimum path such that
\begin{equation}\label{eq-forest}
 (n-1)!\le\alpha b+\beta,
 \qquad
 \wt(P)\ge(n+1)(n-1)!-3+b.
\end{equation}
The forest summation that gives this statement is proved
in~\cite{returns}. The minimum-path decomposition and its local geometric
inputs are supplied by~\cite{framework}. Consequently,
\begin{equation}\label{eq-word-bound}
 S(n)\ge n+(n+1)(n-1)!-3+
 \ceil{\frac{(n-1)!-\beta}{n-1-\delta}}.
\end{equation}
The hereditary condition in~\eqref{eq-affine} is part of the hypothesis;
it cannot be replaced by an inequality for one chosen segment.

\section{Two finite exclusions}
\subsection{The model and its completeness theorem}
Halfar's model records the entry permutation of each block. A piece is
a nonempty list of at most $n-1$ such entries. Consecutive entries in a
piece have an overlap of $n-2$ symbols from the exit of the preceding
block, and consecutive pieces have an overlap of $n-3$ symbols at their
seam. All rotation classes in the entire list are distinct. These
conditions give a relaxation of actual light segments. The deficits
are again $n-1$ minus the piece lengths.

A \emph{window} begins and ends with positive deficit, and each maximal
cluster of positive deficits contains an entry at least two. Write $r$
for the number of pieces and $D$ for the sum of their deficits. We use
positive integers $q,c,b_0$ satisfying $c\le q(n-5)$. The finite exclusion
proves
\begin{equation}\label{eq-window}
 q\bigl(2r+2(n-4)\bigr)+cD
 \le q(n-3)D+b_0
\end{equation}
for every model window, of arbitrary length.

Here is the search and the reason that its finite output is enough.
Relabelling makes the first entry the identity permutation. At a seam,
the next entry is the suffix of length $n-3$ of the preceding exit,
followed by a permutation of its first three symbols. There are six
possibilities. For each, the search tries each piece length from $1$ to
$n-1$; the weight-two rule then determines all entries of that piece.
A branch is rejected if a rotation class has already appeared.

Set
\[
 A=q(n-3)-c,\qquad P=AD,\qquad N=2qr.
\]
The state consists of the last exit, the used rotation classes, $P$, $N$,
and a flag with three possible values. The flag records that the last
piece is full, that its positive cluster consists only of ones, or that
its positive cluster contains an entry at least two. A violation of
\eqref{eq-window} is detected by
\[
 P+b_0<N+2q(n-4)
\]
when the flag permits a window endpoint.

Two further exclusions make the search practical. A full piece cannot
close a cluster consisting only of ones. Also, after a full piece the
search need only start a new positive cluster when $P<N$. In a shortest
counterexample, a failure of the latter condition would leave a shorter
violating suffix. A prefix ending with positive deficit is also discarded
when $P-N\ge2q(n-3)$. A separate finite check excludes a block followed
by $n-2$ full pieces. Thus a following string of full pieces has length
at most $n-3$, while another positive deficit contributes at least
$A-2q\ge0$ to $P-N$; these facts justify the surplus cutoff in a
shortest counterexample.
Halfar's completeness theorem proves these pruning statements together,
by assuming a shortest violating window. It proves that a successful
finite exclusion implies~\eqref{eq-window} for all lengths. Its numerical
encoding and completeness theorem apply for $5\le n\le15$, which includes
both alphabet sizes used here.

The search has an explicit recursion allowance. Exhausting it returns
failure to certify the exclusion, not success. Thus neither a time
limit nor an unexamined branch can establish~\eqref{eq-window}.

\subsection{The capacity calculation}
The actual run sequence also satisfies three local restrictions: a string
of full runs has length at most $n-2$, this drops to $n-3$ when it is
adjacent to a nonfull run, and every consecutive deficit pattern
$(0,1,d)$ has $d\ge2$. These are consequences of the local geometry
in~\cite{framework}. Halfar's list-capacity lemma combines these
restrictions with the window inequality on every factor, accounting
for the two endpoints. An endpoint is charged at most $2q(n-2)+b_0$. For every factor of
an actual light segment, the result is
\[
 4qr\le2AD+2\bigl(2q(n-2)+b_0\bigr).
\]
Thus~\eqref{eq-affine} holds with
\begin{equation}\label{eq-window-parameters}
 \delta=\frac{2q}{A},\qquad
 \beta=\frac{2q(n-2)+b_0}{A}.
\end{equation}
The passage from an actual segment to the model preserves its entries,
rotation classes, and deficits. Applying it to each consecutive factor
supplies exactly the hereditary hypothesis required by the forest
argument. This avoids an assumption that an arbitrary word is already
in the form searched by the program.

The parameters of our two certificates are
\[
\begin{array}{c|rrrr}
 n&c&b_0&q&A\\ \hline
 10&31&508&127&858\\
 13&86&2100&525&5164.
\end{array}
\]
The corresponding recursion allowances are $614$ and $5399$. Both
certificates include the separate exclusion of a block followed by
$n-2$ full pieces. Substitution into~\eqref{eq-window-parameters} gives
\[
\begin{array}{c|ccc}
 n&\delta&\beta&((n-1)!-\beta)/(n-1-\delta)\\ \hline
 10&127/429&1270/429&77837125/1867\\
 13&525/2582&6825/2582&412260708125/10153.
\end{array}
\]
The ceilings in the last column are $41{,}692$ and $40{,}604{,}818$.
Equation~\eqref{eq-word-bound} proves the first two claims of
Theorem~\ref{thm-main}.

The ten-symbol exclusion is evaluated by Lean's kernel, in smaller
subtree certificates. For thirteen symbols, the search is partitioned
into 128 buckets at depth 512, arranged in 32 groups. Some groups are
split further, for a total of 51 native computations. The bucket index
is a rolling hash of the successive piece exits along a search branch. Injectivity is unnecessary:
every bucket is checked, and violations or exhausted recursion before
the partition depth are retained. A kernel proof establishes that the
checks cover the reference search. The additional trust in this result
is confined to the evaluations of these 51 computations.

\section{Why equality is impossible on fourteen symbols}
The remaining improvement comes from the shape of a path that would
attain $B(14)$. The argument has three steps. The forest estimate first
produces one light segment covering all rotation classes. The deficit inequalities
then force every run to have deficit zero or two, with full runs on
both sides of each deficit-two run. Finally, the last symbol of the run
entry is constant. This places all runs in a single partition of the
rotation classes, where a deficit-two run cannot be completed by any
other run.

\subsection{The slack in the forest estimate}
For $n=14$, the parameters from~\cite{returns} are
\[
 p=36,\qquad t=\frac{569}{15},\qquad
 \delta=\frac{1138}{6229},\qquad
 \beta=\frac{15816}{6229},\qquad \alpha=13-\delta.
\]
Suppose, for a contradiction, that a minimum word has length
$93{,}891{,}141{,}149$. Choose a reduced minimum overlap path $P$.
Its weight is $93{,}891{,}141{,}135$, and its available budget is
\[
 b_* = \wt(P)-15\cdot13!+3=485{,}829{,}138.
\]
The margin in~\eqref{eq-forest} is exceptionally small:
\begin{equation}\label{eq-small-margin}
 \alpha b_*+\beta-13!=\frac{1398}{6229}<\frac12.
\end{equation}

We retain two kinds of slack that the usual forest summation discards.
For any
\[
 0<\varepsilon\le\min(1-\delta,\alpha-\beta),
\]
the presence of a component other than the root strengthens the first
inequality in~\eqref{eq-forest} by $\varepsilon$. To see this, mark one
such component in the forest summation. An unabsorbed component costs
at least $\alpha-\beta$. An absorbed component costs at least
$1-\delta$: in the special absorption this is the unused part of its
unit charge, while in the other absorption the two boundary deficits
are at least two and their sum is at least $n$. Substitution into the
endpoint charge $\max(0,d-\delta)$ gives the same surplus. The marked
charge therefore survives summation. This is the forest argument
of~\cite{returns} with one additional charge retained.

The exact weight accounting also has a nonnegative integral residual
$R$, with
\[
 \wt(P)=15\cdot13!-3+b+R.
\]
If $R\ge1$, its contribution to the capacity margin is at least
$\alpha$. Taking $\varepsilon=1/2$ in~\eqref{eq-small-margin} therefore
forces a single component and $R=0$. The components of the forest
partition all rotation classes, so this one component is complete. It is
block-aligned by construction, and its budget is exactly $b_*$. We use
this component for the rest of the argument.

The component has $13!$ blocks, so it has at least two runs.
There is still a possible cost from joins heavier than three and from
nonfull endpoint runs. If $H$ is the total excess of the run-boundary
weights over three, and the first and last deficits are $d_1,d_r$, the
component inequality retains
\[
 \max(0,d_1-\delta)+\max(0,d_r-\delta)+(\alpha-\beta)H.
\]
Each nonzero contribution is at least $1/2$. The same margin excludes
all three. We have obtained one complete light segment with full
endpoint runs, with
\begin{equation}\label{eq-rD}
 r=485{,}829{,}138,\qquad
 D=13r-13!=88{,}757{,}994.
\end{equation}
This step is what allows us to use complete coverage at the end of the
proof. It would be false to assume such a segment at the outset.

\subsection{The deficit sequence}
Write the deficit sequence as positive clusters separated by zeros.
Let $q$ be the number of clusters, $e$ the number of clusters other
than the singleton $[2]$, and $L,R$ the numbers of initial and final
zeros. For each internal zero gap $u$, put
\[
 S_0=\sum_u\max(10-u,0),\qquad M=\#\{u:u=11\}.
\]
The local restrictions of~\cite{returns} say that each cluster $K$ has
$\sum K\ge |K|+1$, that $u\le11$, and that $L+R\le22$. A gap of
length 11 between clusters $K$ and $J$ satisfies the stronger inequality
\begin{equation}\label{eq-critical-gap}
 11+|K|+|J|\le\sum K+\sum J.
\end{equation}
Here $|K|$ denotes the number of entries in the cluster.

Define its reserve by
\[
 \Phi(K)=\frac{6229}{15}(\sum K-2)
       -\frac{1138}{15}(|K|-1)-148\ind_{K\ne[2]},
 \qquad T=\sum_K\Phi(K).
\]
The reserve is nonnegative. If $K\ne[2]$, its mass is at least three,
and direct substitution gives
\begin{equation}\label{eq-exceptional-reserve}
 \Phi(K)\ge\frac{957}{5}.
\end{equation}
For example, among the smallest possible exceptional clusters, the
least reserve occurs at $[1,2]$ or $[2,1]$. Larger mass only increases
the bound. At a gap of length 11, equation~\eqref{eq-critical-gap} gives
\[
 \Phi(K)+\Phi(J)\ge\frac{17207}{5}.
\]
Each cluster is incident with at most two gaps, so
\begin{equation}\label{eq-degree-two}
 \frac{17207}{5}M\le2T.
\end{equation}

The chain-packing inequality of~\cite{returns}, with $p=36$ and
$t=569/15$, gives
\begin{equation}\label{eq-cluster-count}
 4q-148e-\frac{1138}{15}S_0\le144.
\end{equation}
This is its positive-cluster form: regular clusters contribute four;
exceptional clusters and short gaps receive the indicated charges.
It applies to every suffix as well as to the whole sequence. We shall
retain that suffix information shortly.

The total scalar slack is
\[
 E=\frac{6229}{15}D-\frac{1138}{15}r+\frac{15816}{15}
   =\frac{1398}{15},
\]
where the last equality follows from~\eqref{eq-rD}. Counting the entries
in clusters and gaps gives the exact identity
\begin{equation}\label{eq-energy}
\begin{split}
 E={}&144+148e+\frac{1138}{15}S_0-4q+T
       -\frac{1138}{15}M\\
    &+\frac{1138}{15}(22-L-R).
\end{split}
\end{equation}
The first group is nonnegative by~\eqref{eq-cluster-count}, as is the
last term. Equation~\eqref{eq-degree-two} implies
$(1138/15)M\le T/2$. Hence $E\ge T/2$. If an exceptional cluster
existed, equation~\eqref{eq-exceptional-reserve} would give
\[
 E\ge\frac{957}{10}>\frac{1398}{15},
\]
a contradiction. Every positive cluster is therefore $[2]$.
Equation~\eqref{eq-critical-gap} now excludes gaps of length 11.

We need a little more than this classification: each deficit-two run
must have at least two full runs on either side. For $u\le10$, define
\[
 c(u)=
 \begin{cases}
 \dfrac{1138}{15}(10-u)-144,&u\le8,\\
 0,&u=9\text{ or }10.
 \end{cases}
\]
All these costs are nonnegative. The suffix form of the packing
inequality strengthens~\eqref{eq-cluster-count} to
\begin{equation}\label{eq-cut-cost}
 \sum_u c(u)+\frac{1138}{15}(22-L-R)\le E.
\end{equation}
For completeness, this follows by a short induction. In each suffix,
retain the initial clusters whose intervening gaps have length 9 or 10,
and stop at the first shorter gap. Let $v$ be the number of retained
clusters minus $569/30$ times the number of 9-gaps among them. The
chain-packing inequality gives $v\le36$. For one cluster, $v=1$.
Prepending a cluster changes $v$ to $1+v$ across a
10-gap, to $1+v-569/30$ across a 9-gap, and to $1$ across a gap at most
8. In the last case we use the independent suffix bound $v\le36$.
Induction gives
\[
 4q-\frac{1138}{15}S_0+\sum_u c(u)\le4v\le144.
\]
Substitute this inequality into~\eqref{eq-energy}, with $e=T=M=0$, to
obtain~\eqref{eq-cut-cost}.

Since $L,R\le11$, an endpoint with at most nine zeros would contribute
at least $2(1138/15)>E$. Thus $L,R\ge10$. A gap $u\le6$ would have
$c(u)\ge2392/15>E$, so all internal gaps are at least seven. In
particular, every deficit-two run has the required full runs on both
sides.

\subsection{A partition that cannot be filled}
The \emph{pivot} of a run is the last symbol of its entry permutation.
The weight-two rule shows that it is the last symbol of every block
entry in the run. We use the following local consequence of the
six possible weight-three joins.

\begin{lemma}\label{lem-pivot}
Adjacent full runs in a light segment have the same pivot. In a
consecutive triple of runs with deficits $(2,0,0)$, the first two runs
have the same pivot. The reversed assertion holds for $(0,0,2)$.
\end{lemma}
\begin{proof}
For adjacent full runs, the chart-forcing lemma of~\cite{framework}
forces the next entry to be the full-run map of the previous entry.
This map preserves the last letter. For $(2,0,0)$, the full-prefix
forcing lemma applies with prefix length two, equal to the preceding
deficit. Of the six seam entries, it forces the rotation target, whose
last letter is the pivot of the short run. These are respectively the
full-run and full-prefix forms of the same local collision calculation.
For the third assertion, reverse the segment and every block traversal.
A block entered at $(x_1,\ldots,x_{13},s)$ was exited at
$(s,x_1,\ldots,x_{13})$, so its reversed traversal enters at
$(x_{13},\ldots,x_1,s)$. The pivot is unchanged, and the second
assertion applies to the reversed triple. The precise local statements and their
coordinate calculations are included in~\cite{framework}.
\end{proof}

Our deficit sequence consists of zeros and isolated twos, and each two
has at least two zeros on either side. Lemma~\ref{lem-pivot} therefore
shows that all runs have a common pivot $s$.

Fix this symbol $s$. Deleting $s$ from a rotation class leaves a cyclic
ordering of the other thirteen symbols. Conversely, inserting $s$ in
each of the thirteen gaps of a fixed cyclic ordering gives thirteen
rotation classes. Call this set an \emph{insertion loop}. These loops
partition all rotation classes: the deleted cyclic ordering uniquely
determines the loop. For example, deleting $6$ and retaining the cyclic
ordering $12345$ gives a loop of five classes, one for each insertion
gap of $6$.

Each run with pivot $s$ occupies classes from one such loop. A full
run occupies all thirteen classes, and a deficit-two run occupies
eleven. Distinct runs occupy disjoint classes. No two runs can therefore
belong to the same insertion loop, since each occupies more than half
of it. Complete coverage forces each run to fill its loop. Consequently
every run is full, which contradicts the positive value of $D$
in~\eqref{eq-rD}. Equality at $B(14)$ is impossible. Since lengths are
integers, $S(14)\ge B(14)+1$, as required.

\section{The formal statements and their assumptions}
The Palomar version of the accompanying repository submits the ten- and
fourteen-symbol bounds. Its independent statement file defines coverage
directly: every list of $n$ distinct letters from \texttt{Fin n} must
occur as a contiguous factor of the word. It imposes no restriction on
repeated occurrences, the order of first occurrences, or the decomposition
of the word.

The entry point, relative to the repository root, is
\nolinkurl{Submissions/Challenge.lean}. It imports only Mathlib and
also states the general bounds and their asymptotic consequences. The
corresponding \nolinkurl{Submissions/Solution.lean} environment proves
all selected declarations. The file \nolinkurl{AxiomCheck.lean} reports the
assumptions. Both bounds use only \texttt{propext},
\nolinkurl{Classical.choice}, and \nolinkurl{Quot.sound}.

The thirteen-symbol result is retained in this paper for context, but
its native-qualified proof is not included in the Palomar repository
or submitted for registration. The original frozen results repository
at revision \nolinkurl{c9ccb6a0d9ddd1007fdce6efab50807ec47464b9}
contains its statement, solution, and the 51 native computation
assumptions, with their exact types, in
\nolinkurl{lean/tools/native_axioms.json}. Its all-length search
completeness, partition coverage, and word reduction do not depend on
an unproved search heuristic. The 51 computations were also evaluated
with Lean 4.33.0; that original complete proof package is pinned to
Lean 4.34.1. The additional evaluation is a consistency check and does
not remove the native computation assumptions.

The ten-symbol proof is assembled in \nolinkurl{proofs/H20/K10.lean}.
The fourteen-symbol proof is assembled in
\nolinkurl{proofs/new_results/Bounds.lean}; its four main ingredients are
\texttt{ForestRigidity}, \texttt{ScalarRigidity}, \texttt{PivotRigidity},
and \texttt{GlobalPartition}. These are all included. The original
thirteen-symbol module, \nolinkurl{lean/proofs/H20Native/K13.lean},
remains in the frozen repository only.

\begin{thebibliography}{9}
\small
\raggedright
\bibitem{returns}
J. Pantone, \emph{Chain returns and lower bounds for superpermutations},
companion paper in \nolinkurl{papers/lower/}.

\bibitem{framework}
J. Pantone, \emph{The lower-bound framework for superpermutations},
companion Lean library, namespace \texttt{SuperpermutationLowerBound}.
See Sections 1--10 and the refinement modules for the minimum-path,
forest, and local geometry lemmas.

\bibitem{halfar}
J. Halfar, \emph{Three lower bounds in Lean: Theorems A, B and C},
\href{https://github.com/jhalfar/superpermutations/tree/cc5c763bd5854317b20dadb138990802656e6f03/proofs/lean/lower}{lower-bound proof source},
revision \texttt{cc5c763bd5854317b20dadb138990802656e6f03}, October 2026.
The search, completeness, and list-capacity arguments used here are
included with attribution in \nolinkurl{proofs/LowerBounds/}.
\bibitem{halfar-announcement}
J. Halfar, \emph{New upper bound with 1771/3456 and new lower bounds
for $n=8$ to $14$, both in Lean}, announcement to the Superpermutators
group, 6 October 2026.
\url{https://groups.google.com/g/superpermutators/c/6ozhk3faHDA}.
\end{thebibliography}
\end{document}
